  \section{Discussion}
  \label{sec:discussion}

  The results answer the research question for one twin. This section draws out what they mean
  beyond it, and where they stop.

  \subsection{When to code and when to delegate}
  \label{sec:build-or-ask}

  Part (a) has a clear answer for this twin. Delegated with the improved skills, the model reached
  the reference on every inverse question and on N1, as the coded search did, and again in all but
  one of nine repeats of each inverse question, the exception a run of I4 that read the window as
  the original skills did. It did so too in every run of two questions written after the skills,
  one a search over two inputs (\Cref{sec:results-unseen}). With the
  original skills, every answer it gave had the
  reference's verdict, but was only as fine as the values it chose to try
  (\Cref{tab:results-closeness}). No answer, with either version, invented a value where none
  exists or the tools cannot give one, but without skills N2 did (\Cref{sec:results-provenance}). The improved skills called each searched value their own,
  but in the same or the next sentence in only two of seven answers and in none of the six unseen
  runs, and left unlabelled two numbers in I2 that no tool produced.

  Part (b) turns first on what each way must know in advance. Coding needs the question foreseen:
  the coded search has an entry for each, fixing the input to vary, the range, the passing verdicts
  and the deciding budget, and each method assumes a shape for the answer, such as a single change
  from failing to passing in I5. Each unseen question needed a new entry but no new method, since
  both combine the existing ones. Delegating needs only the kind of question foreseen, in one
  section per skill that names no question. Without that section the model either held back, under
  the original skills, or with no skill searched but went past the tools (\Cref{sec:freedom}), so
  this much is needed.

  Time was of the same order: the coded search, running one simulation at a time, was faster on six
  of the seven questions, by a factor of 1.2 to 3.6, and slower on I5, where the model ran
  simulations in parallel (\Cref{tab:results-cost}). Delegating costs tokens instead, on every
  answer: 194 to 473 thousand input tokens with the improved skills, US\$0.25 to 0.42 at list
  prices, against the 130 to 138 thousand, US\$0.16 to 0.20, of a control question,
  which a coded search reached through a skill would roughly match. They follow the model's turns
  rather than its simulations, since a script runs many in one turn: I4 ran 124 simulations for
  401 thousand tokens, I5 ran 15 for 473 thousand. The search over two inputs cost the most: U1
  took up to 2321 simulations, 70 minutes and US\$0.59 per answer, against the coded search's 733
  simulations and 36 minutes. Delegating also gives up what code does by rule. The coded search
  reports a guarded answer beside the raw one, where the improved skills went to the raw limit,
  clearing the budget on I1 and I2 by 0.3 s and on U1 and U2 by 3.0 s against a resolution of
  3.1 s, and it leaves the model no search step to label.

  This suggests a rule. A new or rare kind of question is best delegated, with its kind described
  in the skill, since it is then answered when asked, close to the optimum, with no code written
  for it. A search is worth coding when the question recurs, since a delegated answer pays its
  tokens each time, when the answer must clear the resolution by a rule, or when it must be the
  same every time by construction rather than by observation. I4's runs show the difference:
  they named different seconds as good as the best, and one of ten read the question
  differently and missed it. The heater-boost advisory is the
  coded case: its setpoint search was foreseen, and its guard is part of the tool
  (\Cref{sec:boost}).

  The choice need not be all or nothing. The coded search's entries formulate each question, and
  its three methods are generic. The model formulated every question as those entries do, since
  every scored improved answer reached the reference, and fell short only where code is strong, at the
  guard and the label. The work could be split as Guan et al.\ split planning, where the model
  writes a formal model of the domain and a sound planner searches it \cite{guan2023worldmodels}:
  the model formulates the question when it is asked, a coded method searches, and the twin judges
  each candidate.

  A question delegated often could also be promoted to a coded tool, starting from the model's own
  script. What a twin can answer could then grow from use, while its models stay unchanged: the
  questions operators ask show which capabilities are missing, the model answers them first by
  combining the existing tools, and those that recur are coded as new tools over the same models
  once they pass tests against references, as the coded search did here. The foresight that
  coding needs would then come from the questions asked rather than in advance. This links the
  layer to work that places evolution at the core of engineering DTs, with new services among its
  triggers \cite{alskaif2025evolution}, and to the use of LLMs in modelling DTs
  \cite{yang2025llmdtmodeling}. Neither this promotion nor the split above was tested here.

  That the model searched well does not contradict Valmeekam et al., who find that LLMs
  plan poorly \cite{valmeekam2023planbench}. Each search here is over a single input, and the
  model searched by writing scripts that ran the tools, as program-aided models compute in an
  interpreter \cite{gao2023pal}, choosing new values from earlier results as a model used as an
  optimiser does \cite{yang2024opro}. Xia et al.\ showed such a search in a case study
  \cite{xia2024llmsimulation}, and this work measures one against a reference and a coded search.

  \subsection{Freedom and boundaries in the skills}
  \label{sec:freedom}

  The two versions of the skills differ only in one section per skill and a few lines around it
  (\Cref{app:skills}), yet that difference decided most of the closeness results, for 40\% more
  input tokens on the inverse questions. Choosing what a model is given to read, in its
  instructions and everything else in its context, is now called context engineering
  \cite{mei2025contextsurvey,anthropic2025contextengineering}. In a layer made only of skills, it
  decides how much of the model's ability reaches the twin, through four mechanisms seen in the
  runs.

  \textbf{A boundary is read more widely than it was meant.} The original skills state that the
  model computes and judges nothing itself, that the lid-opening advisory does not search for the
  best moment and that the heater-boost advisory does not look for the longest opening, lines meant
  to keep numerical decisions in the tools. The model mostly read them as ruling out a search: it
  asked back before five of the six inverse questions, declined I6 twice, and in N1 told the
  operator that the tools could not search the hour, presenting the skill's limit as the twin's.
  Context outside the skill moved the same line. In a run of I1 set aside because the assistant
  had also loaded the project's own instructions, which asked whether the layer could answer by
  searching over the tools, it searched at once and reached 253 s, against the scored run's 120 s,
  as did a run with only those instructions added. Both are single runs, observations rather than
  results.

  \textbf{A restricted search stops early.} After the fixed reply the original skills searched on
  grids they chose and reported the best value they had tried, I1 stopping at 120 s after seven
  durations. Without skills no objective was stated either, yet the model refined its search
  unasked, to 252 s on I1 on a 3-second grid and to the limit of 245 s on I2. The coarse searches
  therefore more likely came from the original skills' restrictions than from the missing
  objective. The improved skills lift the restrictions and state the objective, the largest or
  smallest passing value over the whole range and to the second, and reached it on every question.

  \textbf{The objective stated is the one obtained.} ``To the second'' put I1, I2, U1 and U2 at
  the raw limit in all of their runs, inside the resolution. The heater-boost tool reports its
  resolution and guards against it, but the lid-opening tool reports none and its skill asks for
  no guard. In the scored run of I2 the model noticed the 0.3 s margin and advised
  staying ``a little under 4 minutes'', a guard of its own that no tool produced and that its
  repeats did not all give. One run of U2 likewise advised stopping at 4 minutes. Without skills, asked for no precision, the model advised 240 s on both
  I1 and I2, below the limits it had found, and so met the guarded reference. An objective given
  to a model must include the accuracy the twin supports, not only the precision the operator
  acts in.

  \textbf{Without a boundary the search runs past the twin.} Without skills the model searched at
  once on every question, so what the improved skills supply is not the search but its boundary.
  Without one, N2 answered with a room temperature computed on states the twin never
  published and called a high-risk verdict a quirk of the model, and I3 led with a limit the
  lid-opening tool does not apply. The original skills' boundary held the search back, no skill
  let it run past the twin, and the improved skills' boundary allowed the one and stopped the
  other.

  These results show where the line in \Cref{fig:architecture} falls. The original skills drew it
  tighter than the figure: the model computed nothing beyond unit conversion and took no numerical
  decision, which is why the heater-boost advisory leaves a shorter opening unsearched
  (\Cref{sec:boost}). The improved skills move one numerical decision into the language layer,
  which value to run next, but not the judgement: apart from guards of the model's own in I2 and
  in a run of U2, and one count of its own simulations in U1, every number and verdict reported
  came from a run. The line worth drawing is therefore the figure's, that the
  model orchestrates and only the twin predicts and judges, or as the improved skills put it,
  ``Your mathematics chooses what to run; only a run judges it.'' This is the arrangement Kambhampati et
  al.\ argue for in planning, an LLM paired with external model-based verifiers
  \cite{kambhampati2024llmmodulo}, here with the twin as the verifier.

  Mylonas et al.\ enforce such a line in code, rendering an answer's numbers from the tools, and
  without that layer only 39.2\% of answers carried numbers the tools supported
  \cite{mylonas2026governance}. Here, in a different setting, the skills' instructions alone kept
  every number in the scored answers that loaded them to a tool output except I2's two and U1's
  count of its own simulations, while O1
  and O2, which loaded no skill, and five of the eight answers without skills left numbers and
  claims of their own unlabelled. Instructions can carry the grounding, but only where they are
  loaded and as far as they are followed, and the label on a searched value was followed in some
  runs of a question and not in others. This argues for enforcing in code whatever must never
  fail.

  \subsection{Applying the approach to another twin}
  \label{sec:generality}

  Of the three layers in \Cref{fig:architecture}, the structure carries over and the content does
  not. The rules that bind each skill, the split of prediction from judgement, and search as
  repeated simulation do not depend on the incubator. The sources do, and so do the tool parts
  that encode them: the open lid as a scaled heat-loss input, the phases and bands of tempeh, and
  the heater limits. A new capability on the same twin costs a skill, a tool and, where needed,
  new domain knowledge, as the heater boost showed by reusing the simulation and the risk
  assessment unchanged and adding only the search (\Cref{sec:boost}). For a different twin, the
  results suggest five steps.

  \begin{enumerate}
    \item \textbf{Check the twin.} Its model must run forward from outside with the disturbance as
      an input, its current state must be published, and there must be limits to judge a
      prediction against. Had the heat-loss coefficient been a fixed parameter, modelling the open
      lid would have meant changing the twin's code, which this design rules out.
    \item \textbf{Wrap one forward tool per disturbance,} taking inputs in the operator's terms and
      returning a verdict with the numbers that decided it, its range and its resolution, which a
      search and its guard need.
    \item \textbf{Write one skill per advisory} that keeps every reported number to a tool output,
      says what only the tools may judge rather than forbidding the model to compute, and
      describes the search over the tool's inputs with the objective in full, any guard included,
      and the label the search needs.
    \item \textbf{Test the layer alone,} in a session that holds nothing else, against references
      computed by the same tools. The set-aside runs (\Cref{sec:freedom}) also hid a defect in
      both skills, a wrong folder to run the tools from, which the clean runs exposed.
    \item \textbf{Delegate first and code what recurs} (\Cref{sec:build-or-ask}).
  \end{enumerate}

  One twin and one added capability show that the structure allows this reuse, not that it holds
  in general.

  \subsection{Limitations}
  \label{sec:limitations}

  The evidence is narrow: one twin, one model at one effort setting in one assistant, nine
  searches, only U1's over two inputs, and one scored run per question. The nine repeats of each
  inverse question were compared on their answers and labels only, and the original skills' handling of inverse questions changed with what else the
  session held (\Cref{sec:freedom}), so their results describe the layer alone. Larger searches,
  or other models, may differ, and with a slower simulation the number of runs, 124 by the model
  on I4 against the coded search's 39 and up to 2321 on U1, would dominate the time.

  Both sides of the comparison had some hindsight. The improved skills were written after the
  original answers had been read, and their line on phase changes bears on I5
  (\Cref{sec:improved}). The coded search was designed with the reference known, though its rules
  were fixed before it ran. Only the original skills predate the question set, but on two
  questions written after the improved skills, those skills, unchanged, reached the reference in
  every run, as did a coded search designed without it (\Cref{sec:results-unseen}).

  The costs are specific to the set-up. Some improved runs simulated in parallel and the coded
  search did not, and all times depend on the machine, so only their order is compared. The input
  tokens include the assistant's own instructions, sent with every turn, and would differ in
  another assistant.

  The author alone scored the answers, under rules fixed before the runs, with the numbers traced
  by a script and the claims checked by hand.
